\begin{center}\begin{tikzpicture}[font=\small,>=stealth]
\draw[black!30,rounded corners=2pt] (0,0) rectangle (14.4,8.2);
\node[anchor=west,font=\small\bfseries] at (.5,7.72) {例 3：用 $PF=PH$ 把两段距离拉成一条水平线};
\begin{scope}[shift={(4.2,2.35)},scale=.9]
\draw[violet,thick] (-1,-.5)--(-1,4.6) node[above] {$\ell:x=-1$};
\draw[blue,very thick,domain=-.5:4.5,samples=121,variable=\t] plot ({\t*\t/4},\t);
\draw[black!45] (-1.5,0)--(5.5,0);
\draw[orange!85!black,very thick] (-1,2)--(3,2);
\draw[violet,dashed,thick] (3,2)--(4,4)--(-1,4);
\draw[black!45,dashed] (1,0)--(1,2);
\fill (1,0) circle (2pt) node[below] {$F$};
\fill (1,2) circle (2pt) node[above left=6pt] {$P(1,2)$};
\fill (3,2) circle (2pt) node[right] {$A(3,2)$};
\fill (-1,2) circle (2pt) node[left] {$H$};
\fill (4,4) circle (2pt) node[below right=8pt] {$Q$};
\fill (-1,4) circle (2pt) node[left] {$H_Q$};
\end{scope}

\node[anchor=west] at (.7,.55) {$AP+PF=AP+PH\ge d(A,\ell)=4$；$P=(1,2)$ 时等号成立};
\end{tikzpicture}\end{center}
